\subsection{Task 4.1. Bàn giao sản phẩm và phân loại phạm vi}
\label{sec:t4-delivery}

\subsubsection{Sản phẩm bàn giao}

\begin{table}[H]
\centering\small
\caption{Danh mục sản phẩm bàn giao}\label{tab:t4-deliverables}
\begin{tabularx}{\linewidth}{@{}p{.28\linewidth}X@{}}
\toprule
\textbf{Sản phẩm} & \textbf{Nội dung / vị trí}\\\midrule
Phần mềm đang chạy & \url{https://bk-events.vercel.app} (frontend Vercel, backend Render, PostgreSQL Neon, email Brevo).\\
Mã nguồn & \url{https://github.com/giakhang1704/campus_events}, nhánh chính; báo cáo trên nhánh riêng của báo cáo.\\
Hướng dẫn chạy & Hướng dẫn chạy local bằng Docker Compose; hướng dẫn triển khai cloud, vận hành, sao lưu và khôi phục.\\
Tài liệu kỹ thuật & Đặc tả API của từng increment; tài liệu kiến trúc; nhật ký quyết định phạm vi.\\
Bằng chứng chất lượng & Biên bản kết quả kiểm thử I1--I3 và rà soát hệ thống 25/11; bộ test tự động; kiểm thử tự động mỗi lần đẩy mã (CI).\\
Hồ sơ quản lý & Báo cáo này: Baseline V1.0 (Task 2), Baseline V2.0 đề xuất, weekly log, sổ rủi ro, change log, ma trận truy vết.\\
\bottomrule
\end{tabularx}
\end{table}

\subsubsection{Phân loại phạm vi}

\begingroup\small
\renewcommand{\arraystretch}{1.15}
\begin{longtable}{@{}p{.15\linewidth}p{.40\linewidth}p{.38\linewidth}@{}}
\caption{Phạm vi hoàn thành, một phần, hoãn và loại bỏ}\label{tab:t4-scope}\\
\toprule
\textbf{Trạng thái} & \textbf{Hạng mục} & \textbf{Căn cứ / quyết định}\\\midrule\endfirsthead
\toprule
\textbf{Trạng thái} & \textbf{Hạng mục} & \textbf{Căn cứ / quyết định}\\\midrule\endhead
\bottomrule\endfoot
\textbf{Completed} & Tài khoản, đăng nhập, đổi mật khẩu, RBAC 3 vai trò, quản trị tài khoản (FR-01..07). & Test I1 và rà soát 25/11 đạt.\\
 & Tạo, sửa, công bố, tìm và xem sự kiện; giới hạn khoa/năm (FR-08..11). & Test I1, I2 đạt.\\
 & Đăng ký, hủy, danh sách chờ FIFO, 3 loại email (FR-12..18). & Test I2 đạt; thư thật gửi được.\\
 & Điểm danh QR, điểm danh thủ công có lịch sử, báo cáo, xuất CSV (FR-19..23). & Test I3 đạt.\\
 & Triển khai cloud qua HTTPS (NFR-12). & Kiểm tra 27/11.\\
\midrule
\textbf{Partially completed} & Kiểm thử nghiệm thu end-to-end trên cloud (A12.1). & Thu gọn theo CR-05: không chạy end-to-end trên cloud, không thử QR trên điện thoại thật.\\
 & Đo chỉ tiêu M2 (hiệu năng). & Đạt ở local; hạn chế: chưa đo tốc độ thực tế trên cloud (SL-13).\\
 & User Guide (A14.1). & Gộp vào hướng dẫn chạy theo CR-05, chưa có tài liệu riêng cho người dùng cuối.\\
 & Phê duyệt của khách hàng cho các CR do nhóm đề xuất (CR-01..04, CR-06..10) và Baseline V2.0. & Chưa thực hiện do CR-05; trình bày tại demo.\\
\midrule
\textbf{Deferred} & Xuất Excel. & CR-02.\\
 & Lọc sự kiện theo thời gian (TC-BRW-API-08). & CR-10; không thuộc chức năng bắt buộc.\\
 & Thử nghiệm người dùng M5 (5 sinh viên, 20 lượt). & CR-05.\\
 & Email nhắc trước sự kiện, xác minh email/SSO, quên mật khẩu, theo dõi thư đã tới hay bị trả về. & Exclusions ở Task 1.4 hoặc ngoài MVP.\\
\midrule
\textbf{Removed} & Sức chứa không giới hạn (Task 1.4: ``sức chứa nếu có giới hạn''). & CR-03.\\
 & Triển khai trên AWS Lightsail (RA02). & CR-04.\\
\end{longtable}
\endgroup

\subsubsection{Kịch bản demo}

Buổi demo đi qua một luồng end-to-end duy nhất thay vì trình bày từng màn hình rời rạc:
\begin{enumerate}
    \item Quản trị viên cấp tài khoản Người tổ chức; Người tổ chức đăng nhập và đổi mật khẩu tạm.
    \item Người tổ chức tạo sự kiện sức chứa 1, giới hạn một khoa, rồi công bố.
    \item Hai sinh viên đăng ký: người thứ nhất nhận chỗ và email xác nhận, người thứ hai vào danh sách chờ.
    \item Người thứ nhất hủy: hệ thống tự đôn người thứ hai và gửi email.
    \item Người tổ chức mở phiên QR; sinh viên quét bằng điện thoại, xem trước rồi xác nhận.
    \item Người tổ chức sửa điểm danh thủ công có lý do; mở báo cáo và xuất CSV.
    \item Trình bày phần quản lý: Baseline V1.0 và V2.0, change log CR-01..10, sổ rủi ro, số liệu kế hoạch và thực tế; xin khách hàng quyết định các CR chưa được duyệt.
\end{enumerate}
Trước demo, nhóm mở trang khoảng 2 phút để backend Render khởi động (NR1). Phương án dự phòng là chạy bản local bằng Docker Compose.
